\ifdefined\tool\else
\documentclass[uplatex, 10pt, a4j, dvipdfmx]{jsarticle}

\usepackage{packages/abstract}
\usepackage{enumitem}
\usepackage[dvipdfmx]{graphicx}
\usepackage{url}
\usepackage[dvipdfmx, hidelinks]{hyperref}
\usepackage{pxjahyper}

\newcommand{\tool}{VITT}
\newcommand{\AbstractStandalone}{}

% make abstract コマンドでの発表要旨のPDF生成に必要な記述項目
\setAbstractTitle{テストケース可視化ツール\tool{}の試作}
\setAbstractPresenter{有馬 温}

\begin{document}
\AbstractHeader
\begin{AbstractBody}
\fi

ソフトウェアの活用範囲の拡大と複雑化に伴い、ソフトウェア品質の重要性が高まっている。
ソフトウェア品質を確保する工程として、ソフトウェアテストがある。
ソフトウェアテストを実施する準備段階として用意するものの1つに、テストケースがある。
テストケースは、実行事前条件、入力値、アクション(適用可能な場合)、期待結果、および実行事後条件のセットであり、
テスト条件に基づいて記述する。
テストケースの品質を確保するためには、テストケースと要求仕様の対応関係(トレーサビリティ)を明確にすることが重要である。
テストケースの書き方として、テンプレートやフォーマットが提供されており、
これらの必要項目を記述することによって、テストケースは作成できる。

しかし、テストケースと要求仕様との対応関係を確認することは、
テキストで記述したテストケースと要求仕様書の内容を人手で比較する作業が必要であり、手間と時間がかかる。

そこで本研究では、テストケースと要求仕様書の対応関係の確認にかかる時間の削減を目的として、
テストケースを可視化するTCVダイアグラム(Test-Case Visualization Diagram)
の自動生成ツール\tool{}(VIsualization of Testcase Tool)を試作する。
TCVダイアグラムとは、本研究で提案する、テストケースの可視化を目的としたダイアグラムである。
\tool{}は、要求仕様書とテストスイートを入力とし、
TCVダイアグラムを出力する。
なお、テストスイートとは、特定のテストサイクルで実行されるテストケースやテスト手順のセットである。

TCVダイアグラムは、以下を可視化する。

\begin{itemize}
    \item テストケースの充足\\
    ある要求仕様を適正にテストするテストケースが存在する状態
    \item テストケースの過不足\\
    以下のテストケースの抜け、テストケースの不備、テストケースの過剰の3つのテストケースの状態の総称
    \begin{itemize}
        \item テストケースの抜け\\
        ある要求仕様をテストするテストケースが存在しない状態
        \item テストケースの不備\\
        ある要求仕様をテストするテストケースは存在するが、テストケースの内容が要求仕様に対して不十分な状態
        \item テストケースの過剰\\
        テストケースは存在するが、テストする要求仕様が存在しない状態
    \end{itemize}
\end{itemize}

適用例を用いて、本研究で試作した\tool{}が、
正しくTCVダイアグラムを出力することを確認した。

\tool{}の有用性を確認するため、\tool{}を使用した場合と使用しなかった場合の評価実験を行った。
評価実験では、意図的に含めたテストケースの過不足箇所を、
被験者が全て発見するまでにかかる時間を測定した。
この結果、組$\alpha$(閏年判定システム)に対して、
テストケースの過不足箇所の発見にかかる時間を平均で4分28秒(約81.5\%)削減できた。
また、組$\beta$(成績付けシステム)に対して、
テストケースの過不足箇所の発見にかかる時間を平均で3分54秒(約80.7\%)削減できた。
また、評価実験において、\tool{}を使用した場合に、テストケースの過不足箇所の発見漏れはなかったが、
\tool{}を使用しなかった場合は、テストケースの過不足箇所の発見漏れがあった。
これらより、\tool{}は、テストケースの過不足箇所の発見に有用であることを確認した。

以上のことから、\tool{}は、テストケースと要求仕様書の対応関係の確認にかかる時間の削減に有用であるといえる。

\ifdefined\AbstractStandalone
\end{AbstractBody}
\end{document}
\fi
